% ============================================================================
%  PSF Shapelet Score and Moment Score for the Rubin Observatory
%  SPIE Astronomical Telescopes + Instrumentation 2026
%
%  NOTE: spie.cls and spiebib.bst are NOT in this repo yet. Download the SPIE
%  author template and place spie.cls + spiebib.bst in this directory before
%  compiling.  See CLAUDE.md.
% ============================================================================

\documentclass[]{spie}

\usepackage{amsmath}
\usepackage{amssymb}
\usepackage{graphicx}
\usepackage{booktabs}
\usepackage{xcolor}
\usepackage[colorlinks=true, allcolors=blue]{hyperref}

\graphicspath{{fig/}}

% ---- Draft comment macros (remove before submission) -----------------------
\newcommand{\tianqing}[1]{\textcolor{teal}{[TZ: #1]}}
\newcommand{\todo}[1]{\textcolor{red}{[TODO: #1]}}

% ---- Notation --------------------------------------------------------------
\newcommand{\sshape}{S_{\rm shapelet}}
\newcommand{\smom}{S_{\rm moment}}
\newcommand{\bpq}{b_{pq}}
\newcommand{\arcsec}{\ensuremath{^{\prime\prime}}}

\title{The Shapelet Score: An Optics-Sensitive PSF Quality Metric for the
Rubin Observatory}

\author{Tianqing Zhang}
\affiliation{SLAC}

\authorinfo{Send correspondence to T. Zhang: \todo{email}}

\begin{document}
\maketitle

\begin{abstract}
\todo{abstract --- write last, once the results are final.}
\end{abstract}

\keywords{Rubin Observatory, LSST, point spread function, shapelets, image quality,
active optics, wavefront sensing}


% ============================================================================
\section{INTRODUCTION}
\label{intro:0}
The NSF-DOE Vera C. Rubin Observatory will carry out the Legacy Survey of Space and Time
(LSST), imaging the southern sky in six bands over ten years with the 3.2-gigapixel LSST
Camera (LSSTCam) on the 8.4\,m Simonyi Survey Telescope \cite{Ivezic2019}. Many of the
survey's headline science cases --- and cosmic shear in particular --- depend on how well
we understand the point spread function (PSF). The PSF delivered by a ground-based
telescope is a convolution of atmospheric turbulence, the optical system, and the detector,
and errors in modelling it propagate directly into galaxy shape measurements and hence into
cosmological parameter constraints \cite{HirataSeljak2003, Jarvis2021, Zhang2023}.
Monitoring the quality of the delivered PSF, night by night and across the focal plane, is
therefore both an operational requirement for the observatory and a prerequisite for the
survey's precision-cosmology goals.

In practice, image quality is almost always summarised by second-moment quantities: the
full width at half maximum (FWHM) and the two components of the PSF ellipticity. These are
natural choices, they are cheap to compute, and they capture the dominant behaviour of the
PSF. They are also, however, dominated by the atmosphere. Ground-layer seeing sets the FWHM
and contributes most of the ellipticity, so a visit taken through good seeing will report
excellent second-moment image quality even when the optical system is misaligned. The
failure modes that the Rubin active optics system (AOS) exists to correct --- coma,
trefoil, and higher-order asymmetric aberrations arising from mirror figure error, hexapod
misalignment, or incomplete AOS convergence \cite{Xin2015} --- are precisely the modes to
which second moments are least sensitive. A PSF can be compact, nearly round, and still be
badly aberrated. We note that the observatory does track wavefront information directly,
through the Zernike coefficients estimated at the four corner wavefront sensors, but those
estimates sample only the corners of the focal plane and are derived from a separate
curvature-sensing analysis rather than from the science images themselves.

This motivates a metric that is built from the science images, is sensitive to exactly the
asymmetric optical modes that second moments miss, and is insensitive to the atmosphere by
construction. In this work we present such a metric. We decompose each PSF into a polar
shapelet basis \cite{Refregier2003, MasseyRefregier2005, BernsteinJarvis2002} --- the
Gauss-Laguerre eigenfunctions of the two-dimensional quantum harmonic oscillator --- and
show, using a simulated population of atmospheric PSFs, that the shapelet coefficients
split cleanly into two groups: those that respond strongly to atmospheric size and
ellipticity variation, and those that do not respond at all. The second group is exactly
the set of modes with odd azimuthal index $m$, which vanish identically for any
point-symmetric profile. We define the \emph{shapelet score} as the fraction of total
shapelet power carried by those odd-$m$ modes, and a companion \emph{moment score} built
from the third-order adaptive moments of the same stars. Both are dimensionless, both are
computed per PSF, and neither requires any external calibration.

We measure these scores for $1{,}468{,}222$ raft-stacked PSFs drawn from $71{,}015$ LSSTCam
visits spanning 2025 April to 2026 May, covering both the Data Preview 2 (DP2) processing
\cite{DP2} and the subsequent nightly-validation processing. We show that the shapelet
score is essentially independent of the PSF size, that it correlates strongly with the
independently measured moment score, that it responds to AOS wavefront Zernike gradients
across the focal plane, and that it isolates nights and visits with visibly aberrated PSFs
that standard FWHM-based metrics do not flag. We therefore propose the shapelet score as an
image-quality metric that is complementary to --- rather than a replacement for --- PSF
size and the AOS-derived metrics already in use at the observatory.

The paper is organized as follows. In Section~\ref{data:0} we describe the LSSTCam data,
the PSF stamp extraction, and the auxiliary moment and wavefront catalogs. In
Section~\ref{method:0} we set out the shapelet formalism, present the atmospheric
simulation that motivates the mode selection, and define both scores. In
Section~\ref{results:0} we present the measured score distributions, validate the scores
visually and against the higher-order moments, and demonstrate their behaviour as a
function of observing conditions and of time. We conclude in Section~\ref{conclusion:0}.


% ============================================================================
\section{DATA}
\label{data:0}
\todo{roadmap opener}

\subsection{LSSTCam DP2 and post-DP2 visits}
\label{data:visits}
\todo{}

\subsection{PSF star stamps}
\label{data:stamps}
\todo{}

\subsection{Higher-order moment catalog}
\label{data:moments}
\todo{}

\subsection{Active-optics wavefront Zernikes}
\label{data:zernike}
\todo{}


% ============================================================================
\section{PSF QUALITY SCORES}
\label{method:0}
\todo{roadmap opener}

\subsection{Shapelet decomposition of the PSF}
\label{method:shapelet}
\todo{}

\subsection{The shapelet score}
\label{method:shapelet_score}
\todo{}

\subsection{The moment score}
\label{method:moment_score}
\todo{}

\subsection{Score tiers}
\label{method:tiers}
\todo{}


% ============================================================================
\section{RESULTS}
\label{results:0}
\todo{roadmap opener}

\subsection{Score distributions}
\label{results:distributions}
\todo{}

\subsection{Visual validation by tier}
\label{results:gallery}
\todo{}

\subsection{Shapelet score versus moment score}
\label{results:score_comparison}
\todo{}

\subsection{Nightly observatory monitoring}
\label{results:monitoring}
\todo{}

\subsection{Dependence on observing conditions}
\label{results:conditions}
\todo{}

\subsection{Link to optical aberration}
\label{results:optics}
\todo{}


% ============================================================================
\section{CONCLUSION}
\label{conclusion:0}
\todo{conclusion}


% ============================================================================
\acknowledgments
\todo{acknowledgments}

\bibliography{main}
\bibliographystyle{spiebib}

\end{document}
